% !TEX root = ../main.tex
\begin{ouverture}[derivation]{3}{Dérivation}{Partie 2 · Analyse}
{Taux de variation\enspace·\enspace Nombre dérivé\enspace·\enspace Tangentes\enspace·\enspace Calculs}
\objectif{Calculer et interpréter un taux de variation.}
\objectif{Comprendre le nombre dérivé et lire une pente.}
\objectif{Déterminer l'équation d'une tangente.}
\objectif{Dériver des sommes, produits et quotients sur leur domaine.}
\prerequis{Équations de droites et coefficient directeur.}
\prerequis{Calcul littéral, puissances et fractions.}
\end{ouverture}
\essentiel
\essSection{1}{D'une sécante à une tangente}
\begin{definition}
Pour $h\neq0$, le taux de variation entre $a$ et $a+h$ est
$\dfrac{f(a+h)-f(a)}h$. Lorsqu'il se rapproche d'un nombre réel quand $h$
se rapproche de 0, ce nombre est le \textbf{nombre dérivé} $f'(a)$.
\end{definition}
\begin{center}\begin{tikzpicture}\begin{axis}[otLivreGraphe,width=108mm,height=62mm,xmin=-0.5,xmax=3.1,ymin=-1,ymax=8,xtick={0,1,2,3},ytick={1,2,4,6,8}]
\addplot[otNavy,thick,domain=-0.5:2.8,samples=90]{x^2};
\addplot[otGreen,thick,domain=-0.2:3]{2*x-1};
\addplot[otOrange,dashed,domain=0:2.6]{3*x-2};
\addplot[only marks,otNavy]coordinates{(1,1)(2,4)};
\node[above left,font=\small]at(axis cs:1,1){$A$};
\node[above left,font=\small]at(axis cs:2,4){$B$};
\node[otEtiquette,text=otGreenDark]at(axis cs:2.6,3.5){$T$};
\end{axis}\end{tikzpicture}\par
\legende{Pour $f(x)=x^2$, la sécante $(AB)$ a pour pente 3 ; la tangente $T$ en $A$ a pour pente 2.}
\end{center}
\begin{theoreme}[Équation de la tangente]
Si $f$ est dérivable en $a$, la tangente au point $A(a;f(a))$ a pour équation
\[y=f'(a)(x-a)+f(a).\]
Pour un petit $h$, on utilise l'approximation $f(a+h)\approx f(a)+hf'(a)$.
\end{theoreme}
\begin{exemple}
Pour $f(x)=x^2$ et $a=1$, $f(1)=1$ et $f'(1)=2$, donc $T:y=2x-1$.
Ainsi $1{,}02^2\approx1+2\times0{,}02=1{,}04$.
\end{exemple}
\clearpage
\essSection{2}{Les dérivées de référence}
{\renewcommand{\arraystretch}{1.7}\begin{tabularx}{\linewidth}{@{}Y Y Y@{}}
\toprule \textbf{Fonction} & \textbf{Dérivée} & \textbf{Domaine de dérivabilité}\\\midrule
$k$ (constante) & $0$ & $\R$\\
$x^n$, entier $n\geq1$ & $nx^{n-1}$ & $\R$\\
$1/x$ & $-1/x^2$ & $\R\setminus\{0\}$\\
$x^n$, entier $n<0$ & $nx^{n-1}$ & $\R\setminus\{0\}$\\
$\sqrt x$ & $1/(2\sqrt x)$ & $]0;+\infty[$\\\bottomrule
\end{tabularx}}
\essSection{3}{Sommes, produits et quotients}
\begin{propriete}
Pour des fonctions $u$ et $v$ dérivables, et une constante $k$ :
\[(ku+v)'=ku'+v',\qquad(uv)'=u'v+uv'.\]
Sur les intervalles où $v$ ne s'annule pas :
\[\left(\frac uv\right)'=\frac{u'v-uv'}{v^2},\qquad
\left(\frac1v\right)'=-\frac{v'}{v^2}.\]
\end{propriete}
\begin{savoirfaire}[Dériver puis simplifier]
Pour $f(x)=\dfrac{x+1}{x-2}$, on exclut $2$, puis
$f'(x)=\dfrac{(x-2)-(x+1)}{(x-2)^2}=\dfrac{-3}{(x-2)^2}$.
Toujours annoncer le domaine avant de dériver.
\end{savoirfaire}
\begin{attention}
$(uv)'$ n'est pas $u'v'$. Une fonction peut être définie en un point sans y
être dérivable : $\sqrt x$ en 0, ou $|x|$ en 0. Les taux à gauche et à droite
de 0 de $|x|$ valent respectivement $-1$ et $1$.
\end{attention}
\exercices

\rubrique{Découverte}
\begin{blocLivre}{ex-3-1}
\begin{exercice}[1]{Calculer une pente moyenne}
Soit $f(x)=x^2-2x$.
\begin{questions}\item Calculer le taux de variation entre 1 et 3, puis entre 1 et 2.
\item Calculer le taux entre 1 et $1+h$ pour $h\neq0$. En déduire $f'(1)$.
\end{questions}
\end{exercice}
\end{blocLivre}
\begin{blocLivre}{ex-3-2}
\begin{exercice}[1]{Lire la dérivée sur une tangente}
La tangente à une courbe en $A(2;3)$ passe aussi par $B(4;7)$.
\begin{questions}\item Déterminer $f(2)$ et $f'(2)$.
\item Écrire l'équation de cette tangente.
\end{questions}
\label{figlivre-r03-02-tangente-ex}
\figureLivre{r03-02-tangente}

\end{exercice}
\end{blocLivre}
\begin{blocLivre}{ex-3-3}
\begin{exercice}[1]{Dérivées immédiates}
Donner les dérivées de $f(x)=5x^4-3x^2+7x-8$, $g(x)=3/x$,
et $h(x)=\sqrt x+2x$. Préciser les domaines de dérivabilité.
Donner aussi la dérivée de $k(x)=x^{-3}$ sur $\R\setminus\{0\}$.
\end{exercice}
\end{blocLivre}
\begin{blocLivre}{ex-3-4}
\begin{exercice}[1]{Une tangente à un polynôme}
$f(x)=x^2+3x-1$.
\begin{questions}\item Calculer $f'(x)$, $f(1)$ et $f'(1)$.
\item Déterminer la tangente au point d'abscisse 1.
\end{questions}
\end{exercice}
\end{blocLivre}
\begin{blocLivre}{ex-3-5}
\begin{exercice}[1]{Reconnaître l'opération}
Sans développer, dériver $u(x)=(2x+1)(x^2-3)$ et
$v(x)=\dfrac{x+1}{x-2}$. Préciser les domaines.
\end{exercice}
\end{blocLivre}
\clearpage
\rubrique{Entraînement}
\begin{blocLivre}{ex-3-6}
\begin{exercice}[2]{Revenir à la définition}
On considère $f(x)=x^2$ et un réel $a$.
\begin{questions}
\item Montrer que $\dfrac{f(a+h)-f(a)}h=2a+h$ pour $h\neq0$.
\item En déduire $f'(a)$, puis la tangente au point d'abscisse $a$.
\item Pour $a=2$, calculer les pentes des sécantes correspondant à $h=1$, $0{,}1$, $-0{,}1$ et $0{,}01$.
\item Pour un pas $h>0$, écrire une fonction Python \code{secantes(h, n)} qui renvoie la liste
des coefficients directeurs des sécantes passant par les points d'abscisses $kh$ et $(k+1)h$,
pour $k$ allant de 0 à $n-1$. Prévoir \code{secantes(0.5, 4)} et comparer à $f'(kh)$.
\end{questions}
\end{exercice}
\end{blocLivre}
\begin{blocLivre}{ex-3-7}
\begin{exercice}[2]{Dériver trois produits}
Calculer et simplifier les dérivées de :
\begin{questions}\item $f(x)=(x^2+1)(3x-2)$ sur $\R$.
\item $g(x)=x^2\sqrt x$ sur $]0;+\infty[$.
\item $h(x)=(x-1)^2(x+2)$ sur $\R$, en développant d'abord ou en utilisant la règle du produit.
\end{questions}
\end{exercice}
\end{blocLivre}
\begin{blocLivre}{ex-3-8}
\begin{exercice}[2]{Dériver trois quotients}
Préciser le domaine de dérivabilité et calculer la dérivée de :
\begin{tasks}(3)
\task $f(x)=\dfrac{2x-1}{x+3}$.
\task $g(x)=\dfrac{x^2+1}{x}$.
\task $h(x)=\dfrac{x}{x^2+1}$.
\end{tasks}
Expliquer pourquoi une simplification préalable peut être utile pour $g$.
\end{exercice}
\end{blocLivre}
\begin{blocLivre}{ex-3-9}
\begin{exercice}[2]{Des tangentes parallèles}
$f(x)=x^3-3x+1$.
\begin{questions}\item Calculer $f'(x)$.
\item Déterminer les abscisses où la tangente est parallèle à $y=9x-2$.
\item Donner les équations des deux tangentes et vérifier qu'elles sont distinctes.
\end{questions}
\end{exercice}
\end{blocLivre}
\begin{blocLivre}{ex-3-10}
\begin{exercice}[2]{Approcher une racine carrée}
Soit $f(x)=\sqrt x$.
\begin{questions}\item Déterminer la tangente à sa courbe en $x=4$.
\item Utiliser cette tangente pour approcher $\sqrt{4{,}08}$.
\item Comparer à une valeur donnée par la calculatrice et calculer l'erreur absolue à $10^{-6}$ près.
\end{questions}
\end{exercice}
\end{blocLivre}
\begin{blocLivre}{ex-3-11}
\begin{exercice}[2]{Deux exemples de non-dérivabilité}
\begin{questions}
\item Pour $f(x)=|x|$, calculer $\dfrac{f(h)-f(0)}h$ lorsque $h>0$, puis $h<0$.
Conclure sur la dérivabilité en 0.
\item Pour $g(x)=\sqrt x$ et $h>0$, montrer que le taux entre 0 et $h$ vaut $1/\sqrt h$.
Pourquoi cela ne donne-t-il pas un nombre dérivé réel en 0 ?
\item Ces fonctions sont-elles malgré tout définies en 0 ?
\end{questions}
\label{figlivre-r03-11-nonderivable-ex}
\figureLivre{r03-11-nonderivable}

\end{exercice}
\end{blocLivre}
\clearpage
\rubrique{Problèmes}
\begin{blocLivre}{ex-3-12}
\begin{exercice}[3]{Tangentes à une parabole}
On considère $\Cf:y=x^2$ et le point $P(0;-1)$.
\begin{center}\begin{tikzpicture}\begin{axis}[otLivreGraphe,width=97mm,height=56mm,xmin=-2.3,xmax=2.3,ymin=-2,ymax=5,xtick={-2,-1,0,1,2},ytick={-1,1,3,5}]
\addplot[otNavy,thick,domain=-2.2:2.2,samples=80]{x^2};
\addplot[otGreen,domain=-2.3:2.3]{2*x-1};
\addplot[otGreen,domain=-2.3:2.3]{-2*x-1};
\addplot[only marks,otNavy]coordinates{(0,-1)(1,1)(-1,1)};
\node[below right,font=\small]at(axis cs:0,-1){$P$};
\end{axis}\end{tikzpicture}\end{center}
\begin{questions}
\item Écrire l'équation $T_a$ de la tangente au point d'abscisse $a$.
\item Montrer que $P$ appartient à $T_a$ si et seulement si $a^2=1$.
\item Donner les points de contact et les deux équations. Retrouver le dessin.
\end{questions}
\end{exercice}
\end{blocLivre}
\begin{blocLivre}{ex-3-13}
\begin{exercice}[3]{Vitesse moyenne et vitesse instantanée}
Un mobile a pour position $s(t)=t^3-6t^2+12t$, en mètres, pour $t\in[0;4]$ secondes.
\begin{questions}
\item Calculer la vitesse moyenne entre $t=1$ et $t=3$.
\item Calculer la vitesse instantanée $v(t)=s'(t)$ et l'écrire sous forme factorisée.
\item Quand le mobile est-il instantanément à l'arrêt ? Calculer sa position à cet instant.
\item Le mobile peut-il avoir une vitesse négative selon ce modèle ? Justifier.
\end{questions}
\end{exercice}
\end{blocLivre}
\begin{blocLivre}{ex-3-14}
\begin{exercice}[3]{Retrouver un polynôme grâce à une tangente}
On cherche $f(x)=ax^2+bx+c$ telle que $f(0)=1$, $f(2)=3$ et $f'(2)=4$.
\begin{questions}
\item Traduire les trois conditions en équations sur $a$, $b$ et $c$.
\item Déterminer ces coefficients et vérifier les trois conditions.
\item Écrire la tangente $T$ en $x=2$.
\item Factoriser $f(x)-(4x-5)$. La courbe est-elle au-dessus ou au-dessous de $T$ ?
\end{questions}
\end{exercice}
\end{blocLivre}
\begin{blocLivre}{ex-3-15}
\begin{exercice}[3]{Coût marginal et coût supplémentaire}
Le coût de fabrication de $q$ objets est modélisé, pour $q\in[0;100]$, par
$C(q)=0{,}02q^3-0{,}6q^2+10q+100$ euros.
\partieExo{A}{Calculer et interpréter}
\begin{questions}\item Calculer $C'(q)$ et le coût marginal $C'(20)$.
\item Calculer exactement $C(21)-C(20)$ ; comparer à $C'(20)$.
\item À l'aide de l'approximation affine, estimer $C(20{,}1)-C(20)$.
\end{questions}
\partieExo{B}{Contrôler une approximation}
\begin{questions}[resume]
\item Développer $C(20+h)-C(20)$ et montrer que cette différence vaut
$10h+0{,}6h^2+0{,}02h^3$.
\item Pour $h=0{,}1$, calculer l'erreur de l'approximation de la question 3.
Expliquer pourquoi une dérivée ne donne pas toujours exactement le coût d'un objet supplémentaire.
\end{questions}
\end{exercice}
\end{blocLivre}
